\chapter{Montgomery-curve mathematics}

This chapter formalises the arithmetic behind Curve25519 / X25519: the
projective $x$-only differential-addition and doubling formulas, the
Montgomery ladder, the transport of the Montgomery curve to Weierstrass
form (so its group law comes for free from Mathlib), and the capstone
correctness theorem for the RFC~7748 base point. Four facts about
Curve25519 that are expensive to certify inside Lean (primality of the
base field, primality of the subgroup order, the base point and its
order) are bridged across the Rocq$\leftrightarrow$Lean boundary from the
AUCurves development and recorded as axioms; everything downstream of them
is proved.

\section{The \texorpdfstring{$x$}{x}-only formulas}

\begin{definition}[Projective $x$-only doubling]
  \label{def:xdbl}
  \lean{CatCrypt.Crypto.ECC.xdbl}
  \leanok
  For a Montgomery curve constant $a_{24} = (A+2)/4$, the projective
  doubling map on a pair $(X,Z)$ representing $x = X/Z$ (Costello--Smith
  eq.~12) is
  \[
    \mathsf{xdbl}\,a_{24}\,(X,Z) \;=\;
      \bigl(\,(X+Z)^2 (X-Z)^2,\;\; c\,((X-Z)^2 + a_{24}\,c)\,\bigr),
    \qquad c = (X+Z)^2 - (X-Z)^2,
  \]
  where $c = 4XZ$. Its output represents $x(2\cdot P)$ when $(X,Z)$
  represents $x(P)$.
\end{definition}

\begin{definition}[Projective $x$-only differential addition]
  \label{def:xdadd}
  \lean{CatCrypt.Crypto.ECC.xdadd}
  \leanok
  Given projective $x$-coordinates of $P$, $Q$ and their difference
  $P-Q$ (Costello--Smith eq.~13),
  \begin{align*}
    \mathrm{DA} &= (X_P - Z_P)(X_Q + Z_Q), &
    \mathrm{CB} &= (X_P + Z_P)(X_Q - Z_Q), \\
    \mathsf{xdadd}\,(X_P,Z_P)\,(X_Q,Z_Q)\,(X_D,Z_D)
      &= \bigl(\,Z_D\,(\mathrm{DA}+\mathrm{CB})^2,\;\;
                  X_D\,(\mathrm{DA}-\mathrm{CB})^2\,\bigr).
  \end{align*}
  Its output represents $x(P+Q)$.
\end{definition}

\begin{lemma}[Symmetry of $\mathsf{xdadd}$]
  \label{lem:xdadd-symm}
  \uses{def:xdadd}
  \lean{CatCrypt.Crypto.ECC.xdadd_symm}
  \leanok
  $\mathsf{xdadd}\,PZ\,QZ\,DZ = \mathsf{xdadd}\,QZ\,PZ\,DZ$: the difference
  term $\mathrm{DA}-\mathrm{CB}$ is squared, so swapping the two summands
  leaves the result unchanged.
\end{lemma}

\begin{definition}[Montgomery curve parameters]
  \label{def:monty-params}
  \lean{CatCrypt.Crypto.ECC.MontyParams}
  \leanok
  A parameter record over a commutative ring $F$ carrying the curve
  coefficient $A$ (for $y^2 = x^3 + A x^2 + x$) and the derived constant
  $a_{24}$, kept as an independent field subject to the well-formedness
  invariant $4\,a_{24} = A + 2$ (avoiding division in the representation).
\end{definition}

\begin{definition}[Curve25519 parameters]
  \label{def:curve25519-params}
  \uses{def:monty-params}
  \lean{CatCrypt.Crypto.ECC.curve25519Params}
  \leanok
  The Montgomery parameters of Curve25519: $A = 486662$ and
  $a_{24} = 121666 = (486662 + 2)/4$.
\end{definition}

\section{Projective equivalence and homogeneity}

\begin{definition}[Projective equivalence]
  \label{def:proj-equiv}
  \lean{CatCrypt.Crypto.ECC.ProjEquiv}
  \leanok
  Two pairs are $x$-equivalent, written $(X_1,Z_1) \sim (X_2,Z_2)$, iff
  $X_1 Z_2 = X_2 Z_1$; i.e.\ they carry the same $x$-coordinate
  $X/Z$ in the field of fractions. The relation is reflexive and
  symmetric.
\end{definition}

\begin{lemma}[Restricted transitivity of $\sim$]
  \label{lem:proj-equiv-trans}
  \uses{def:proj-equiv}
  \lean{CatCrypt.Crypto.ECC.ProjEquiv.trans}
  \leanok
  Over a field, if $p \sim q$, $q \sim r$, and $q$ is nondegenerate
  ($q.1 \neq 0$ or $q.2 \neq 0$), then $p \sim r$. Transitivity can fail
  when the middle term is $(0,0)$, so nondegeneracy is required.
\end{lemma}

\begin{theorem}[$\mathsf{xdbl}$ respects $\sim$]
  \label{thm:xdbl-projequiv}
  \uses{def:xdbl, def:proj-equiv}
  \lean{CatCrypt.Crypto.ECC.xdbl_projEquiv}
  \leanok
  If $p \sim q$ then $\mathsf{xdbl}\,a_{24}\,p \sim \mathsf{xdbl}\,a_{24}\,q$.
  Both coordinates of $\mathsf{xdbl}$ are homogeneous of degree $4$ in
  $(X,Z)$, so a common scaling of the input scales the output uniformly.
\end{theorem}

\begin{theorem}[$\mathsf{xdadd}$ respects $\sim$]
  \label{thm:xdadd-projequiv}
  \uses{def:xdadd, def:proj-equiv}
  \lean{CatCrypt.Crypto.ECC.xdadd_projEquiv}
  \leanok
  If $p \sim p'$, $q \sim q'$, and $d \sim d'$ then
  $\mathsf{xdadd}\,p\,q\,d \sim \mathsf{xdadd}\,p'\,q'\,d'$: each of
  $\mathrm{DA}+\mathrm{CB} = 2(X_P X_Q - Z_P Z_Q)$ and
  $\mathrm{DA}-\mathrm{CB} = 2(X_P Z_Q - X_Q Z_P)$ is homogeneous in each
  argument, and the $X_D/Z_D$ factors cancel.
\end{theorem}

\section{The abstract correspondence: \texorpdfstring{$\mathsf{MontyCurveGroup}$}{MontyCurveGroup}}

\begin{definition}[$\mathsf{MontyCurveGroup}$]
  \label{def:monty-curve-group}
  \uses{def:xdbl, def:xdadd, def:monty-params}
  \lean{CatCrypt.Crypto.ECC.MontyCurveGroup}
  \leanok
  A hypothesis class packaging the Costello--Smith~\S4 correspondence: an
  abelian group $G$ of curve points, an $x$-coordinate projection
  $\mathsf{xProj} : G \to F \times F$, and the two specifications
  \begin{align*}
    \text{(doubling)}\quad&\mathsf{xdbl}\,a_{24}\,(\mathsf{xProj}\,P)
      \;\sim\; \mathsf{xProj}\,(P+P), \\
    \text{(addition)}\quad&\mathsf{xdadd}\,(\mathsf{xProj}\,P)\,
      (\mathsf{xProj}\,Q)\,(\mathsf{xProj}\,(P-Q))
      \;\sim\; \mathsf{xProj}\,(P+Q),
  \end{align*}
  stated up to projective equivalence.
\end{definition}

\section{The Montgomery ladder (group level)}

\begin{definition}[Bits to natural]
  \label{def:bits-to-nat}
  \lean{CatCrypt.Crypto.ECC.bitsToNat}
  \leanok
  The natural number read from a bit list, MSB first:
  $\mathsf{bitsToNat}(b::bs) = 2\,\mathsf{bitsToNat}(bs) + [b]$ and
  $\mathsf{bitsToNat}\,[] = 0$.
\end{definition}

\begin{definition}[Ladder step]
  \label{def:ladder-step}
  \lean{CatCrypt.Crypto.ECC.ladderStep}
  \leanok
  One double-and-add step on a state $(R_0, R_1)$ of an abelian group,
  branching on bit $b$:
  \[
    \mathsf{ladderStep}\,(R_0,R_1)\,b =
      \begin{cases}
        (R_0 + R_1,\; R_1 + R_1) & b = \mathsf{true},\\
        (R_0 + R_0,\; R_0 + R_1) & b = \mathsf{false}.
      \end{cases}
  \]
\end{definition}

\begin{definition}[Ladder fold]
  \label{def:ladder-fold}
  \uses{def:ladder-step}
  \lean{CatCrypt.Crypto.ECC.ladderFold}
  \leanok
  Iterating $\mathsf{ladderStep}$ over a bit list (head processed first,
  MSB first) from the initial state $(0, P)$.
\end{definition}

\begin{definition}[Ladder invariant]
  \label{def:ladder-inv}
  \lean{CatCrypt.Crypto.ECC.LadderInv}
  \leanok
  The predicate $\mathsf{LadderInv}\,P\,a\,(R_0,R_1)$ asserting
  $R_0 = a \cdot P$ and $R_1 = (a+1)\cdot P$.
\end{definition}

\begin{theorem}[Montgomery ladder correctness (algebraic)]
  \label{thm:ladder-fold-correct}
  \uses{def:ladder-fold, def:ladder-inv, def:bits-to-nat}
  \lean{CatCrypt.Crypto.ECC.ladderFold_correct}
  \leanok
  For any abelian group $G$, any $P : G$ and any bit list $bits$,
  \[
    \mathsf{LadderInv}\,P\,(\mathsf{bitsToNat}\,bits)\,
      (\mathsf{ladderFold}\,P\,bits),
  \]
  i.e.\ the fold ends in state $(n\cdot P, (n+1)\cdot P)$ with
  $n = \mathsf{bitsToNat}\,bits$. Proved by induction on bit length.
\end{theorem}

\section{The field-level \texorpdfstring{$x$}{x}-only ladder}

\begin{definition}[Field-level ladder step]
  \label{def:xladder-step}
  \uses{def:xdbl, def:xdadd}
  \lean{CatCrypt.Crypto.ECC.xladderStep}
  \leanok
  One step of the $x$-only ladder on projective state
  $(R_0, R_1) \in (F\times F)^2$, carrying the base $x$-coordinate $x_P$
  as the invariant difference:
  \[
    \mathsf{xladderStep}\,a_{24}\,x_P\,(R_0,R_1)\,b =
      \begin{cases}
        (\mathsf{xdadd}\,R_0\,R_1\,x_P,\; \mathsf{xdbl}\,a_{24}\,R_1) & b,\\
        (\mathsf{xdbl}\,a_{24}\,R_0,\; \mathsf{xdadd}\,R_0\,R_1\,x_P) & \neg b.
      \end{cases}
  \]
  This is the pure field program a Jasmin/bedrock2/Rust implementation runs.
\end{definition}

\begin{definition}[Field-level ladder fold]
  \label{def:xladder-fold}
  \uses{def:xladder-step}
  \lean{CatCrypt.Crypto.ECC.xladderFold}
  \leanok
  Iterating $\mathsf{xladderStep}$ over a bit list from the initial
  projective state $((1,0),\,x_P)$, where $(1,0)$ represents the point at
  infinity.
\end{definition}

\begin{definition}[Field-level ladder invariant]
  \label{def:xladder-inv}
  \uses{def:proj-equiv, def:monty-curve-group}
  \lean{CatCrypt.Crypto.ECC.XLadderInv}
  \leanok
  The projective invariant tying the field state to the abstract group:
  $\mathsf{XLadderInv}\,P\,n\,(s_0,s_1)$ holds iff
  $s_0 \sim \mathsf{xProj}\,(n\cdot P)$ and
  $s_1 \sim \mathsf{xProj}\,((n+1)\cdot P)$.
\end{definition}

\begin{theorem}[Field-level step preserves the invariant]
  \label{thm:xladder-step-preserves}
  \uses{def:xladder-step, def:xladder-inv, thm:xdbl-projequiv, thm:xdadd-projequiv, lem:proj-equiv-trans, def:monty-curve-group}
  \lean{CatCrypt.Crypto.ECC.xladderStep_preserves_invariant}
  \leanok
  Over a field, given negation-invariance of $\mathsf{xProj}$ and
  nondegeneracy of the intermediate $\mathsf{xdbl}$/$\mathsf{xdadd}$
  outputs, if $\mathsf{XLadderInv}\,P\,n\,s$ holds then so does
  $\mathsf{XLadderInv}\,P\,(2n + [b])\,(\mathsf{xladderStep}\,a_{24}\,
  (\mathsf{xProj}\,P)\,s\,b)$. Composes
  \ref{thm:xdbl-projequiv}, \ref{thm:xdadd-projequiv} and
  \ref{lem:proj-equiv-trans} with the group specs.
\end{theorem}

\section{Montgomery as a Weierstrass curve}

\begin{definition}[Montgomery curve as Weierstrass curve]
  \label{def:montgomeryW}
  \lean{CatCrypt.Crypto.ECC.montgomeryW}
  \leanok
  The equation $y^2 = x^3 + A x^2 + x$ is literally the Weierstrass curve
  with $(a_1,a_2,a_3,a_4,a_6) = (0, A, 0, 1, 0)$. Instantiating Mathlib's
  \texttt{WeierstrassCurve} at these coefficients reuses its proved
  \texttt{AddCommGroup} structure on affine points.
\end{definition}

\begin{lemma}[Weierstrass equation is the Montgomery equation]
  \label{lem:montgomeryW-equation}
  \uses{def:montgomeryW}
  \lean{CatCrypt.Crypto.ECC.montgomeryW_equation}
  \leanok
  A point $(x,y)$ lies on $\mathsf{montgomeryW}\,A$ iff
  $y^2 = x^3 + A x^2 + x$.
\end{lemma}

\begin{definition}[Projective $x$-coordinate map]
  \label{def:xprojW}
  \uses{def:montgomeryW}
  \lean{CatCrypt.Crypto.ECC.xProjW}
  \leanok
  The map sending a Weierstrass affine point to its projective
  $x$-coordinate: the point at infinity $\mapsto (1,0)$, and a finite
  point $(x,y) \mapsto (x,1)$. It is negation-invariant, since Montgomery
  negation flips only $y$.
\end{definition}

\begin{theorem}[$\mathsf{xdbl}$ doubling spec on Weierstrass points]
  \label{thm:xdbl-spec-w}
  \uses{def:xdbl, def:xprojW, def:proj-equiv, lem:montgomeryW-equation}
  \lean{CatCrypt.Crypto.ECC.xdbl_spec_W}
  \leanok
  For $2 \neq 0$ and $4\,a_{24} = A + 2$, and every point $P$ on
  $\mathsf{montgomeryW}\,A$,
  \[
    \mathsf{xdbl}\,a_{24}\,(\mathsf{xProjW}\,A\,P)
      \;\sim\; \mathsf{xProjW}\,A\,(P+P),
  \]
  covering the point at infinity, the $2$-torsion case, and the generic
  case.
\end{theorem}

\begin{theorem}[$\mathsf{xdadd}$ addition spec on Weierstrass points]
  \label{thm:xdadd-spec-w}
  \uses{def:xdadd, def:xprojW, def:proj-equiv, lem:montgomeryW-equation}
  \lean{CatCrypt.Crypto.ECC.xdadd_spec_W}
  \leanok
  For $2 \neq 0$ and all points $P, Q$ on $\mathsf{montgomeryW}\,A$,
  \[
    \mathsf{xdadd}\,(\mathsf{xProjW}\,A\,P)\,(\mathsf{xProjW}\,A\,Q)\,
      (\mathsf{xProjW}\,A\,(P-Q))
      \;\sim\; \mathsf{xProjW}\,A\,(P+Q),
  \]
  covering $P = 0$, $Q = 0$, $P = \pm Q$, and the generic
  ($x_P \neq x_Q$) case. Reduces to a polynomial identity over the field
  using only the two curve equations.
\end{theorem}

\begin{definition}[Weierstrass transport into $\mathsf{MontyCurveGroup}$]
  \label{def:montgomeryW-mcg}
  \uses{def:monty-curve-group, def:montgomeryW, def:xprojW, thm:xdbl-spec-w, thm:xdadd-spec-w}
  \lean{CatCrypt.Crypto.ECC.montgomeryW_MontyCurveGroup}
  \leanok
  Packaging: the Weierstrass group on $(\mathsf{montgomeryW}\,A)$'s affine
  points, with $x$-map $\mathsf{xProjW}\,A$, is a
  $\mathsf{MontyCurveGroup}$ instance whenever $2 \neq 0$ and
  $4\,a_{24} = A + 2$. This discharges the Costello--Smith~\S4
  correspondence with no separate Montgomery-form group-law port.
\end{definition}

\section{Curve25519}

\begin{definition}[Curve25519 base-field prime]
  \label{def:curve25519-prime}
  \lean{CatCrypt.Crypto.ECC.curve25519Prime}
  \leanok
  $p = 2^{255} - 19$.
\end{definition}

\begin{definition}[Curve25519 base field]
  \label{def:fp25519}
  \uses{def:curve25519-prime}
  \lean{CatCrypt.Crypto.ECC.Fp25519}
  \leanok
  The type alias $\mathbb{F}_p = \mathbb{Z}/p\mathbb{Z}$ for
  $p = 2^{255} - 19$.
\end{definition}

\begin{definition}[Base-field primality (AUCurves bridge, taken as an axiom)]
  \label{def:curve25519-prime-prime}
  \uses{def:curve25519-prime}
  \lean{CatCrypt.Crypto.ECC.curve25519Prime_prime}
  \leanok
  $\mathsf{Nat.Prime}\,(2^{255} - 19)$. Proving this in Lean needs a costly
  Pocklington/Pratt certificate; it is discharged in the Rocq AUCurves
  development and imported across the Rocq$\leftrightarrow$Lean boundary as
  an axiom. It unlocks $\mathsf{Field}\,\mathbb{F}_p$.
\end{definition}

\begin{definition}[Concrete Curve25519 point group]
  \label{def:curve25519-point}
  \uses{def:montgomeryW, def:fp25519}
  \lean{CatCrypt.Crypto.ECC.Curve25519Point}
  \leanok
  The affine points of $\mathsf{montgomeryW}\,(486662)$ over
  $\mathbb{F}_p$, an $\mathsf{AddCommGroup}$ inherited from Mathlib.
\end{definition}

\begin{definition}[Curve25519 $\mathsf{MontyCurveGroup}$ instance]
  \label{def:curve25519-mcg}
  \uses{def:curve25519-point, def:montgomeryW-mcg, def:curve25519-params, def:curve25519-prime-prime}
  \lean{CatCrypt.Crypto.ECC.curve25519_MontyCurveGroup}
  \leanok
  Specialising \ref{def:montgomeryW-mcg} to $A = 486662$,
  $a_{24} = 121666$ over $\mathbb{F}_p$ (using $2 \neq 0$ and
  $4\cdot 121666 = 486662 + 2$): the $x$-only formulas correctly compute
  scalar multiplication on $\mathsf{Curve25519Point}$.
\end{definition}

\begin{definition}[X25519 scalar multiplication]
  \label{def:x25519}
  \lean{CatCrypt.Crypto.ECC.x25519}
  \leanok
  $\mathsf{x25519}\,k\,P = k \cdot P$ in the group. For the Curve25519
  base point this is X25519 scalar multiplication.
\end{definition}

\begin{theorem}[Ladder computes X25519]
  \label{thm:x25519-eq-ladder}
  \uses{def:x25519, def:ladder-fold, thm:ladder-fold-correct}
  \lean{CatCrypt.Crypto.ECC.x25519_eq_ladder}
  \leanok
  $\mathsf{x25519}\,(\mathsf{bitsToNat}\,bits)\,P
     = (\mathsf{ladderFold}\,P\,bits).1$, a corollary of
  \ref{thm:ladder-fold-correct}.
\end{theorem}

\section{The subgroup order and the base point (AUCurves bridges)}

\begin{definition}[Curve25519 subgroup order]
  \label{def:curve25519-suborder}
  \lean{CatCrypt.Crypto.ECC.curve25519SubgroupOrder}
  \leanok
  $\ell = 2^{252} + 27742317777372353535851937790883648493$, the order of
  Curve25519's prime-order subgroup and of the RFC~7748 base point.
\end{definition}

\begin{definition}[Subgroup order primality (AUCurves bridge, taken as an axiom)]
  \label{def:curve25519-suborder-prime}
  \uses{def:curve25519-suborder}
  \lean{CatCrypt.Crypto.ECC.curve25519SubgroupOrder_prime}
  \leanok
  $\mathsf{Nat.Prime}\,\ell$. Bridged from the AUCurves development
  ($E/E[4]$ has prime order $\ell$); imported as an axiom exactly as
  \ref{def:curve25519-prime-prime}. This underpins the cyclic-subgroup
  argument that $n\cdot B \neq 0$ for $0 < n < \ell$.
\end{definition}

\begin{definition}[Curve25519 base point (AUCurves bridge, taken as an axiom)]
  \label{def:curve25519-basepoint}
  \uses{def:curve25519-point}
  \lean{CatCrypt.Crypto.ECC.curve25519_basepoint}
  \leanok
  A concrete $\mathsf{Curve25519Point}$ standing for the RFC~7748 base
  point $u = 9$. Bridged from AUCurves and taken as an axiom; only its
  additive order is constrained (see \ref{def:curve25519-basepoint-addorder}),
  so downstream results hold for any order-$\ell$ point.
\end{definition}

\begin{definition}[Base-point order (AUCurves bridge, taken as an axiom)]
  \label{def:curve25519-basepoint-addorder}
  \uses{def:curve25519-basepoint, def:curve25519-suborder}
  \lean{CatCrypt.Crypto.ECC.curve25519_basepoint_addOrder}
  \leanok
  $\mathsf{addOrderOf}\,B = \ell$: the base point has additive order equal
  to the subgroup order ($E\_basepoint\_order$, AUCurves). Bridged and
  taken as an axiom.
\end{definition}

\section{Ladder nondegeneracy and the capstone}

\begin{definition}[Ladder nondegeneracy]
  \label{def:ladder-nondeg}
  \uses{def:xdbl, def:xdadd, def:curve25519-point, def:curve25519-mcg}
  \lean{CatCrypt.Crypto.ECC.LadderNondeg}
  \leanok
  The predicate asserting that, at every scalar $n \le \mathsf{bitsToNat}\,bits$
  along the ladder path over $P$, the intermediate $\mathsf{xdbl}$ and
  $\mathsf{xdadd}$ outputs are projectively nondegenerate (not $(0,0)$).
  This is the side condition needed to invoke
  \ref{lem:proj-equiv-trans} at each step.
\end{definition}

\begin{theorem}[Field-level ladder correctness for Curve25519]
  \label{thm:xladder-curve25519-correct}
  \uses{def:xladder-fold, def:xladder-inv, thm:xladder-step-preserves, def:curve25519-mcg, def:ladder-nondeg}
  \lean{CatCrypt.Crypto.ECC.xladderFold_curve25519_correct}
  \leanok
  For a point $P$, a bit list, and $\mathsf{LadderNondeg}\,P\,bits$, the
  field-level fold $\mathsf{xladderFold}$ satisfies $\mathsf{XLadderInv}$
  with the abstract-group ladder. Proved by induction, stepping with
  \ref{thm:xladder-step-preserves}.
\end{theorem}

\begin{theorem}[End-to-end X25519 correctness, conditional on nondegeneracy]
  \label{thm:x25519-ladder-correct}
  \uses{thm:xladder-curve25519-correct, def:x25519, def:ladder-nondeg}
  \lean{CatCrypt.Crypto.ECC.x25519_ladder_correct}
  \leanok
  Given $\mathsf{LadderNondeg}\,P\,bits$, the first component of the
  field-level ladder run on $\mathsf{xProj}\,P$ is projectively equivalent
  to $\mathsf{xProj}\,(\mathsf{x25519}\,(\mathsf{bitsToNat}\,bits)\,P)$ ---
  correctness wherever the ladder path stays nondegenerate.
\end{theorem}

\begin{theorem}[Base point discharges nondegeneracy]
  \label{thm:curve25519-basepoint-laddernondeg}
  \uses{def:ladder-nondeg, def:curve25519-basepoint, def:curve25519-basepoint-addorder, def:curve25519-suborder-prime}
  \lean{CatCrypt.Crypto.ECC.curve25519_basepoint_LadderNondeg}
  \leanok
  For any bit list whose value is below the subgroup order $\ell$, the
  entire ladder path over the base point $B$ is projectively nondegenerate:
  $\mathsf{LadderNondeg}\,B\,bits$. Since $\ell$ is an odd prime, no
  non-identity multiple $n\cdot B$ ($0 < n < \ell$) is a $4$-torsion point,
  which is exactly what forces each $\mathsf{xdbl}$/$\mathsf{xdadd}$ output
  to be nondegenerate.
\end{theorem}

\begin{theorem}[Unconditional X25519 ladder correctness for the base point]
  \label{thm:x25519-ladder-correct-basepoint}
  \uses{thm:x25519-ladder-correct, thm:curve25519-basepoint-laddernondeg, def:curve25519-prime-prime, def:curve25519-suborder-prime, def:curve25519-basepoint, def:curve25519-basepoint-addorder}
  \lean{CatCrypt.Crypto.ECC.x25519_ladder_correct_basepoint}
  \leanok
  \textbf{Capstone.} For every scalar below the subgroup order $\ell$
  (i.e.\ $\mathsf{bitsToNat}\,bits < \ell$), the first component of the
  field-level Montgomery ladder run on $\mathsf{xProj}\,B$ is projectively
  equivalent to $\mathsf{xProj}\,(\mathsf{x25519}\,(\mathsf{bitsToNat}\,bits)\,B)$:
  \[
    (\mathsf{xladderFold}\,a_{24}\,(\mathsf{xProj}\,B)\,bits).1
      \;\sim\; \mathsf{xProj}\,(\mathsf{x25519}\,(\mathsf{bitsToNat}\,bits)\,B).
  \]
  This is \ref{thm:x25519-ladder-correct} with the $\mathsf{LadderNondeg}$
  side condition \emph{discharged} via
  \ref{thm:curve25519-basepoint-laddernondeg}: it carries \emph{no}
  nondegeneracy hypothesis and is unconditional on the RFC~7748 scalar
  range $k < \ell$. It rests only on the four AUCurves-bridged axioms
  (\ref{def:curve25519-prime-prime}, \ref{def:curve25519-suborder-prime},
  \ref{def:curve25519-basepoint}, \ref{def:curve25519-basepoint-addorder}).
\end{theorem}
